\documentclass[letterpaper, 12pt]{article}
\input{tex/_preamble}
\begin{document}
\flushleft\includegraphics[width=0.5\textwidth]{img/home/241017_final_logo_mockup.png}

\cosigsection{Formulaic research}
{7 May 2026}

Many research fields have seen recent explosions in the number of articles performing what could be described as ``formulaic'' research. Such articles often follow a similar template and similar methods and are only differentiated by the exact variable, outcome, disease, subject, etc. that has been swapped in. This strain of research has been described as ``cut-and-paste research'', \href{https://doi.org/10.1371/journal.pmed.1004851}{fast-churn research} and \href{https://doi.org/10.1126/science.zgawnij}{``research Mad Libs\textregistered''}. 

This guide describes several fields, methodologies and datasets that have recently seen a large number of formulaic research articles published. Many are suspected to result from the activity of research paper mills, organizations that are known to sell authorship on mass-produced manuscripts.

When evaluating any manuscript, some helpful questions to ask include:

\begin{itemize}
    \setlength\itemsep{-0.5em}
    \item Is the hypothesis well-justified?
    \item Are the methodological specifics (e.g., inclusion/exclusion criteria) well-justified?
    \item Does this work seem to resemble a previous work with only minor variations?
    \item Does this work control for possible false discoveries (see COSIG's \href{https://osf.io/csxd5}{entry on multiple hypothesis correction})?
    \item Is the data under analysis trustworthy?
\end{itemize}

\subsection*{Example 1: Mendelian randomization studies}

\href{https://doi.org/10.1038/s43586-021-00092-5}{Mendelian randomization (MR)} is an epidemiological method that leverages genotype information in large-scale cohorts (e.g., the \href{https://www.ukbiobank.ac.uk/}{UK Biobank}) to study the causal relationship between an exposure (e.g., physical activity, drinking coffee, etc.) and an outcome (e.g., Alzheimer's disease, type II diabetes, etc.). \href{https://doi.org/10.1093/ije/dyx028}{Two-sample MR (2SMR)} is a variant of MR that relies solely on summary statistics, typically from \href{https://doi.org/10.1038/s43586-021-00056-9}{genome-wide associations studies (GWAS)}. Recently, thousands of poor-quality \href{https://doi.org/10.1093/ije/dyx028}{two-sample MR (2SMR)} studies have been published, likely aided by the accessibility of large genetic databases and the availability of user-friendly interfaces for interacting with this data (e.g., the R packages \href{https://doi.org/10.7554/eLife.34408}{TwoSampleMR} and \href{https://doi.org/10.12688/wellcomeopenres.19995.2}{MendelianRandomization}). These studies often only report a 2SMR study with no additional supporting evidence for an association and feature weakly-motivated hypothetical associations (such as an association between \href{http://dx.doi.org/10.3389/fendo.2024.1396032}{air pollution and non-alcoholic fatty liver disease} or \href{http://dx.doi.org/10.3390/nu15092091}{noodle consumption and metabolic syndrome}).

\subsubsection*{Additional resources}

\begin{itemize}
    \setlength\itemsep{-0.5em}
    \item \href{https://doi.org/10.1186/s12944-024-02284-w}{``Reclaiming mendelian randomization from the deluge of papers and misleading findings'' (2024)}
    \item \href{https://doi.org/10.1136/egastro-2025-100187}{``Attention to the misuse of Mendelian randomisation in medical research'' (2025)}
    \item \href{https://doi.org/10.1016/S2213-8587(23)00348-0}{``Mendelian randomisation at 20 years: how can it avoid hubris, while achieving more?'' (2024)}
    \item \href{https://doi.org/10.1007/s10654-025-01317-7}{``The rapid growth in Mendelian randomization studies'' (2025)}
\end{itemize}

\subsection*{Example 2: Associative research using National Health and Nutrition Examination Survey (NHANES) data}

The \href{https://www.cdc.gov/nchs/nhanes/index.html}{National Health and Nutrition Examination Survey (NHANES)} is a survey of thousands of United States adults and children conducted annually by the Centers for Disease Control and Prevention (CDC), who have made the data freely available since 1999. The data includes surveys, health exams and laboratory tests. Thousands of studies have been published analyzing NHANES data. However, recent years have seen a proliferation of studies performing formulaic single-factor analysis using NHANES data, often on unjustifiably small slices of the surveyed population and only over a limited number of years. For instance, \href{https://doi.org/10.3389/fendo.2023.1245199}{Nie et al. (2023)} analyze a hypothetical association between systematic immune-inflammation index and diabetes but only over the four years 2017 to 2020 while suitable data was available for the 21 years between 1999 and 2020.

Many of these studies report on specious hypotheses with little statistical support. These studies also recycle a small number of statistical methods, especially logistical regression, use of cubic splines, or further subdividing populations into quartiles, in order to cross statistical thresholds. For instance, \href{https://doi.org/10.1186/s12888-023-04935-1}{Zhang et al. (2023)} report on an association between higher serum albumin concentration and depressive symptoms with an odds ratio of 0.98 (95\% confidence interval 0.95-0.99) and a p-value of 0.037.

\subsubsection*{Additional resources}

\begin{itemize}
    \setlength\itemsep{-0.5em}
    \item \href{https://doi.org/10.1371/journal.pbio.3003152}{``Explosion of formulaic research articles, including inappropriate study designs and false discoveries, based on the NHANES US national health database'' (2025)}
    \item \href{https://doi.org/10.3897/ese.2025.e165043}{``Meeting the challenges posed by mass-produced manuscripts and click-data science'' (2025)}
    \item \href{https://doi.org/10.1101/2025.09.09.25335401}{``Dramatic increases in redundant publications in the Generative AI era'' (2025)}
    \item \href{https://doi.org/10.1126/science.z6eukgo}{``Journals and publishers crack down on research from open health data sets'' (2025)}
\end{itemize}

\subsection*{Example 3: Pharmacovigilance signal detection studies using FDA Adverse Event Reporting System (FAERS) data}

The \href{https://www.fda.gov/drugs/fdas-adverse-event-reporting-system-faers/fda-adverse-event-reporting-system-faers-public-dashboard}{FDA Adverse Events Reporting System (FAERS)} database contains reports of \href{https://en.wikipedia.org/wiki/Adverse_event}{adverse events} (termed Individual Case Safety Reports, ICSRs), medication errors and product quality problems submitted to the \href{https://fda.gov}{United States Food and Drug Administration (FDA)} for human drugs. As of 2025, FAERS includes more than 35 million reports, originating mainly from the pharmaceutical industry, patients and physicians. Pharmacovigilance researchers can leverage this data to suggest new causal associations between an intervention (e.g., taking a specific pain-relieving medication, taking a combination of medications) and an adverse event (e.g., heart palpitations, weight loss).

Given the size of the data, pharmaceutical companies, regulatory agencies and pharmacovigilance researchers typically rely on statistical data mining algorithms to aid this process. The most common method for data mining is \href{https://doi.org/10.1007/s40264-024-01421-9}{disproportionality analysis}, which involves identifying drug-event combinations that occur more frequently than expected by chance. Statistics reported in disproportionality analyses include the Proportional Reporting Ratio (PRR), Reporting Odds Ratio (ROR) and the Empirical Bayes Geometric Mean (EBGM, an output of the Multi-item Gamma Poisson Shrinker algorithm, MGPS).

ICSRs and disproportionality analyses using them are subject to inherent limitations:

\begin{itemize}
    \setlength\itemsep{-0.5em}
    \item Reports are made voluntarily and various external factors can influence the likelihood of a report actually being submitted to FAERS.
    \item Drug-event associations are not certain and there is often insufficient information included in ICSRs to rule out other factors involved each case.
    \item FAERS does not provide denominators, e.g., the number of patients exposed to a drug or experiencing an adverse event. Note that the total number of adverse events of any kind associated with a drug is often used as a proxy for exposure in disproportionality analyses.
    \item Disproportionality measures are subject to \href{https://doi.org/10.1007/s40264-013-0063-5}{competition bias}---an over-reporting of cases for a particular event or drug can artificially lower disproportionality measures for other events and drugs.
\end{itemize}

Disproportionality methods do not allow researchers to estimate incidence and infer causality. This prevents researchers from ranking drugs according to their risk, especially when comparing medications with heterogeneous clinical uses and time on the market. As a hypothesis-generating study design, disproportionality analysis can only detect the so-called ``signal of disproportionate reporting'' (SDR). The magnitude of this signal does not necessarily correlate with the strength of the association or its likelihood of representing a real causal relationship.

Since the mid-2000s, the FDA has provided public access to FAERS data through the \href{https://fis.fda.gov/sense/app/95239e26-e0be-42d9-a960-9a5f7f1c25ee/sheet/7a47a261-d58b-4203-a8aa-6d3021737452/state/analysis}{FAERS Public Dashboard}, the \href{https://open.fda.gov/}{OpenFDA API} and \href{https://www.fda.gov/drugs/fdas-adverse-event-reporting-system-faers/fda-adverse-event-reporting-system-faers-latest-quarterly-data-files}{quarterly data reports}. In the last few years, FAERS data has frequently been used to publish formulaic research articles featuring disproportionality analyses. These articles are generally not well-motivated, apply multiple generic methods (leading to redundant results) and interpret results in an nonspecific, standardized format overlooking inherent and study-specific limitations.

In one such study, \href{https://doi.org/10.1111/andr.13533}{Wang et al. (2023)} perform a disproportionality analysis focused on identifying adverse events associated with sildenafil, a widely-used drug for which adverse events are largely well-understood. They apply four different algorithms to compute SDRs even though these algorithms are heavily inter-dependent and render very similar results to one another. The authors find that sildenafil has an ``unexpected'' association with pulmonary hypertension. However, sildenafil is frequently prescribed to treat pulmonary hypertension. Other events for which the authors find an association include malignant melanoma and aortic dissection rupture, conditions which are more common in the older patient population that takes sildenafil. The authors do not adjust for age in their analysis.

These articles often use similar formatting and style (same figures, methods and tables) and contain obvious errors in technical terminology. For example, disproportionality analysis methods like ``Bayesian confidence propagation neural network'' are frequently reported as ``Bayesian confidence interval progressive neural network'', ``Bayesian belief propagation neural network'' and ``Bayesian configuration promotion neural network''. 

\subsubsection*{Vaccine Adverse Events Reporting System (VAERS) data}

The limitations described above for analyzing FAERS data are also applicable to the analysis of data from the \href{https://vaers.hhs.gov/}{Vaccine Averse Events Reporting System (VAERS)}. VAERS enables anyone to report any adverse event they believe to be related to a vaccine they received. VAERS allows public health officials to detect adverse reactions to vaccines that are so rare that even large, pre-approval clinical trials are underpowered to detect a causal association before a vaccine is approved and marketed. For instance, VAERS reports prompted the FDA and Centers for Disease Control and Prevention (CDC) to investigate \href{https://archive.cdc.gov/#/details?url=https://www.cdc.gov/vaccines/vpd-vac/rotavirus/vac-rotashield-historical.htm}{an association between intestinal intussusception and the rotavirus vaccine RotaShield\textsuperscript{\textregistered} in infants}. The resulting investigations concluded that RotaShield\textsuperscript{\textregistered} increased the risk for intussusception by 1 or 2 cases per 10,000 vaccinated infants. RotaShield\textsuperscript{\textregistered} was withdrawn from the market by its manufacturer in 1999.

VAERS has frequently been misused to make broad claims about vaccine safety that are not supported by evidence. For more information, see reports from \href{https://publichealth.jhu.edu/2022/what-vaers-is-and-isnt}{John Hopkins University Bloomberg School of Public Health} and \href{https://www.mcgill.ca/oss/article/covid-19-critical-thinking-health/dont-fall-vaers-scare-tactic}{McGill University Office for Science and Society}.

\subsubsection*{Additional resources}

\begin{itemize}
    \setlength\itemsep{-0.5em}
    \item \href{https://doi.org/10.5281/zenodo.17116884}{``The Rising Misuse of Pharmacovigilance Reporting Systems: A Threat to Evidence-Based Medicine'' (2025)}
    \item \href{https://doi.org/10.1002/cpt.3701}{``FDA Adverse Event Reporting System (FAERS) Essentials: A Guide to Understanding, Applying, and Interpreting Adverse Event Data Reported to FAERS'' (2025)}
    \item \href{https://doi.org/10.3389/fdsfr.2023.1323057}{``Conducting and interpreting disproportionality analyses derived from spontaneous reporting systems'' (2024)}
    \item \href{https://doi.org/10.1007/s40264-024-01423-7}{``The REporting of A Disproportionality Analysis for DrUg Safety Signal Detection Using Individual Case Safety Reports in PharmacoVigilance (READUS-PV): Explanation and Elaboration'' (2024)}
    \item \href{https://doi.org/10.1007/s40264-024-01421-9}{``The Reporting of a Disproportionality Analysis for Drug Safety Signal Detection Using Individual Case Safety Reports in PharmacoVigilance (READUS-PV): Development and Statement.'' (2024)}
    \item \href{https://www.mcgill.ca/oss/article/covid-19-critical-thinking-health/dont-fall-vaers-scare-tactic}{``Don’t Fall for the `VAERS Scare' Tactic'' (2021)}
    \item \href{https://publichealth.jhu.edu/2022/what-vaers-is-and-isnt}{``What VAERS Is (And Isn’t)'' (2022)}
\end{itemize}


\subsection*{Example 4: Bibliometric analyses}

\href{https://en.wikipedia.org/wiki/Bibliometrics}{Bibliometrics/scientometrics} is the quantitative study of indicators of knowledge transfer and development. One such indicator of knowledge transfer is citations between scientific articles. The number of citations an article has received is often used as a heuristic of that article's importance or quality and is a common analytical variable in bibliometric analyses. While high-quality \href{https://en.wikipedia.org/wiki/Metascience}{metascientific} research routinely uses such heuristics, the accessibility of bibliometric data through services like \href{https://scopus.com}{Scopus} and \href{/https://webofscience.com/}{Web of Science} makes it an easy target for low-quality, formulaic research articles. 

For instance, thousands of articles have been published where the authors report on characteristics of the N highest-cited article in a particular field (e.g., \href{https://doi.org/10.1097/coa.0000000000000021}{``A Bibliometric Analysis of the 100 Most Cited Articles in Cornea'', 2023}, \href{https://doi.org/10.3389/fphar.2022.963032}{``A bibliometric analysis of the 100 most-cited articles on curcumin'', 2022}). Such articles often report summary statistics of these top-cited articles (such as describing the distribution of countries of origin), but more often than not without comparison to a control group or to all other articles in the field, giving these quantitative analyses little analytical value.

Formulaic metascience articles can also report on author characteristics in any selected field of the scientific literature. For instance, \href{https://doi.org/10.7759/cureus.47208}{Djahanshahi et al. (2023)} report on gender trends among first authors in the scientific literature on \href{https://en.wikipedia.org/wiki/Wolff%E2%80%93Parkinson%E2%80%93White\_syndrome}{Wolff–Parkinson–White syndrome}, a rare heart disease for which fewer than 200 articles have been written. The authors' reasoning for selecting such a narrow field for analysis is unclear.

\subsubsection*{Additional resources}

\begin{itemize}
    \setlength\itemsep{-0.5em}
    \item \href{https://reeserichardson.blog/2024/12/30/recent-encounters-with-atom-thin-salami-slicing/}{``Recent encounters with atom-thin salami slicing: Case \#2: Same paper mill, same procedure, different subject, different authors'' (2024)}
\end{itemize}

\subsection*{Example 5: Querying a large language model}

Since the public release of \href{https://en.wikipedia.org/wiki/ChatGPT}{ChatGPT} in 2022, thousands of \href{https://scholar.google.com/scholar?start=10&q="conversation+with+chatgpt"+OR+"interview+with+chatgpt"&hl=en}{``conversations'' and ``interviews'' with ChatGPT} have been published in peer-reviewed and non-peer-reviewed outlets. These commentary pieces can be written with little effort on practically any scientific topic. Many are severely limited in scope, typically only sharing a single chain of queries and outputs on a single topic, and do little to interrogate the veracity of the large language model's responses to the user's queries. Moreover, because the companies that develop these models publicly release new versions and quietly make updates between official versions, these conversations are typically not reproducible and any insights gained from them quickly obsolesce.

\subsection*{Example 6: Non-coding RNAs}

Many suspected paper mill products, especially those in the field of \href{https://en.wikipedia.org/wiki/Non-coding_RNA}{non-coding RNAs} and disease, resemble variations on a theme. Plausible-sounding papers can be generated by combining just about any non-coding gene with any human disease and any molecular pathway. This pattern can even be observed in article titles alone, as noted for a collection of papers dubbed the \href{https://scienceintegritydigest.com/2020/02/21/the-tadpole-paper-mill/}{``Tadpole paper mill''}.

\begin{table}[h!tbp]
\scriptsize
\centering
\begin{tabular}{ccccccc}
  \textbf{\begin{tabular}[c]{@{}c@{}}Insert the name\\of a molecule\end{tabular}} &
  \textbf{\begin{tabular}[c]{@{}c@{}}Pick a verb\\(present tense,\\third person\\singular form)\end{tabular}} &
  \textbf{\begin{tabular}[c]{@{}c@{}}Choose\\one or two\\cellular\\processes\end{tabular}} &
  \textbf{\begin{tabular}[c]{@{}c@{}}Pick a\\cancer or\\cell type\end{tabular}} &
  \textbf{\begin{tabular}[c]{@{}c@{}}Pick your\\connector\\word\end{tabular}} &
  \textbf{\begin{tabular}[c]{@{}c@{}}Choose a verb\\(present\\particle form)\end{tabular}} &
  \textbf{\begin{tabular}[c]{@{}c@{}}Insert name\\of miRNA\\or pathway\end{tabular}} \\ \\ \hline
  \multicolumn{1}{|c|}{\begin{tabular}[c]{@{}c@{}}<protein name>\\<drug name>\\<RNA name>\end{tabular}} &
  \multicolumn{1}{c|}{\begin{tabular}[c]{@{}c@{}}alleviates\\attenuates\\exerts\\governs\\inhibits\\prevents\\promotes\\protects\\relieves\\remits\\retards\\suppresses\end{tabular}} &
  \multicolumn{1}{c|}{\begin{tabular}[c]{@{}c@{}}apoptosis\\autophagy\\inflammation\\invasion\\migration\\proliferation\\ viability\end{tabular}} &
  \multicolumn{1}{c|}{\begin{tabular}[c]{@{}c@{}}lung cancer\\medulloblastoma\\renal carcinoma\\ovarian cancer\end{tabular}} &
  \multicolumn{1}{c|}{\begin{tabular}[c]{@{}c@{}}by\\via\\through\end{tabular}} &
  \multicolumn{1}{c|}{\begin{tabular}[c]{@{}c@{}}activating\\attenuating\\declining\\downregulating\\inhibiting\\interfering with\\modulating\\targeting\\regulating\\upregulating\end{tabular}} &
  \multicolumn{1}{c|}{\begin{tabular}[c]{@{}c@{}}<miRNA name>\\<pathway name>\\<protein name>\end{tabular}} \\ \hline
\end{tabular}
    \caption*{Template structure for article titles from the Tadpole paper mill. Adapted from a \href{https://scienceintegritydigest.com/2020/02/21/the-tadpole-paper-mill/}{2020 blog post by Elisabeth Bik.}}
\end{table}

These articles often feature strikingly similar templates, images and data re-used from other articles to refer to different experiments, non-functional reagents (see COSIG's \href{https://osf.io/2egvz}{entry on the subject}), experiments in misidentified and non-verifiable cell lines (see COSIG's \href{https://osf.io/d7we5}{entry on the subject}) and poorly-motivated hypotheses.

\subsubsection*{Additional resources}

\begin{itemize}
    \setlength\itemsep{-0.5em}
    \item \href{https://doi.org/10.1002/1873-3468.13201}{``Systematic fabrication of scientific images revealed'' (2018)}
    \item \href{https://doi.org/10.1177/1177271919829162}{``The Possibility of Systematic Research Fraud Targeting Under-Studied Human Genes: Causes, Consequences, and Potential Solutions'' (2019)}
    \item \href{https://doi.org/10.1002/1873-3468.13747}{``Digital magic, or the dark arts of the 21st century—how can journals and peer reviewers detect manuscripts and publications from paper mills?'' (2021)}
    \item \href{https://scienceintegritydigest.com/2020/02/21/the-tadpole-paper-mill/}{Science Integrity Digest: ``The Tadpole Paper Mill'' (2020)}
    \item \href{https://doi.org/10.1093/nar/gkac1139}{``Protection of the human gene research literature from contract cheating organizations known as research paper mills'' (2022)}

\end{itemize}

\subsection*{Example 7: Metaphor-based metaheuristics}

\href{https://en.wikipedia.org/wiki/Metaheuristic}{Metaheuristics} are conceptual frameworks for developing optimization and search algorithms. Many metaheuristics are inspired by natural phenomena and processes. For instance, one of the most popularly-applied metaheuristic approaches is \href{https://en.wikipedia.org/wiki/Simulated_annealing}{simulated annealing}, a stochastic algorithm inspired by \href{https://en.wikipedia.org/wiki/Annealing_(materials_science)}{metallurgical annealing}. Other nature-inspired metaheuristics that have proven to be both popular and useful include \href{https://en.wikipedia.org/wiki/Ant_colony_optimization_algorithms}{ant colony optimization} and \href{https://en.wikipedia.org/wiki/Genetic_algorithm}{genetic algorithms}.

Following on the popularity of these metaphor-based metaheuristics, recent years have seen thousands of articles published proposing new metaphor-based metaheuristics, particularly optimization algorithms based on the behavior of animals. Many of these proposed methods are poorly-described and are of dubious generalizability and novelty. Often, the algorithm is highly similar to a well-established metaheuristic, just couched in language derived from a new metaphor. For instance, the popular \href{https://doi.org/10.1177/003754970107600201}{harmony search algorithm}, inspired by how jazz bands improvise, \href{https://doi.org/10.4018/978-1-4666-0270-0.ch005}{has been shown} to be equivalent to a special case of the well-established metaheuristic \href{https://en.wikipedia.org/wiki/Evolution_strategy}{evolutionary strategies}.

The \href{https://fcampelo.github.io/EC-Bestiary/}{Evolutionary Computation Bestiary} has collected hundreds of articles proposing ``novel'' metaphor-based metaheuristics. These metaphors include the mating behavior of \href{https://doi.org/10.1007/978-981-13-3708-6_18}{barnacles}, \href{https://doi.org/10.1088/1757-899x/782/5/052028}{beetles} and \href{https://doi.org/10.1007/s00521-019-04464-7}{naked mole rats}, the growth of \href{https://doi.org/10.1016/j.asoc.2015.07.045}{tumors} and the organization of \href{http://dx.doi.org/10.4236/ijis.2014.41002}{soccer leagues}, even concepts like \href{https://doi.org/10.1109/CINTI.2010.5672231}{reincarnation} and \href{https://doi.org/10.1016/j.isatra.2014.03.018}{interior design}. Many such metaheuristics are trivial variations on the popular \href{https://en.wikipedia.org/wiki/Particle_swarm_optimization}{particle swarm optimization} method. As stated by \href{https://doi.org/10.1111/itor.12001}{S\"orensen (2013)}:

\begin{quote}
    \textit{...``novel'' metaheuristics based on new metaphors should be avoided if they cannot demonstrate a contribution to the field. To stress the point: renaming existing concepts does not count as a contribution. Even though methods may be called ``novel'' by their originator, many present no new ideas, except for the occasional marginal variant of an already existing method. Moreover, these methods take up the space of truly innovative ideas and research, for example in the analysis of existing heuristics. Because these methods invariably change the vocabulary, they are difficult to understand. Combined with the fact that the authors of these methods usually neglect to properly position ``their'' method in the metaheuristics literature, such methods present a loss of time and a step backward rather than forward.}
\end{quote}

\subsubsection*{Additional resources}

\begin{itemize}
    \setlength\itemsep{-0.5em}
    \item \href{https://doi.org/10.48550/arXiv.1704.00853}{``A history of metaheuristics'' (2017)}
    \item \href{https://doi.org/10.1111/itor.12001}{``Metaheuristics—the metaphor exposed'' (2013)}
    \item \href{https://doi.org/10.1007/978-3-030-60376-2_10}{``Grey Wolf, Firefly and Bat Algorithms: Three Widespread Algorithms that Do Not Contain Any Novelty'' (2020)}
    \item \href{https://doi.org/10.4018/978-1-4666-0270-0.ch005}{``A Rigorous Analysis of the Harmony Search Algorithm: How the Research Community can be Misled by a “Novel” Methodology'' (2012)}
    \item \href{https://doi.org/10.1111/itor.13176}{``Exposing the grey wolf, moth-flame, whale, firefly, bat, and antlion algorithms: six misleading optimization techniques inspired by bestial metaphors'' (2022)}
    \item \href{https://fcampelo.github.io/EC-Bestiary/}{The Evolutionary Computation Bestiary}
\end{itemize}

\subsection*{Example 8: Kaggle datasets}

\href{https://www.kaggle.com/datasets}{Kaggle} is a platform where users can share any and all kinds of data. Many Kaggle datasets are uploaded as a part of \href{https://www.kaggle.com/competitions}{public competitions}, such as a competition where users compete to create the best algorithm for \href{https://www.kaggle.com/competitions/ieee-fraud-detection/overview}{detecting fraud in e-commerce transactions}. Users often publicly share their entry's code on the Kaggle platform.

A growing number of articles analyze datasets sourced from Kaggle. Many of these articles propose machine learning approaches to tasks on these datasets that do not markedly differ from nor improve upon the performance achieved by publicly-available solutions on a competition's ``Leaderboard'' page. For instance, articles by \href{https://doi.org/10.1007/s44163-024-00169-6}{de Filippis and Al Foysal (2024)} and \href{https://doi.org/10.1038/s41598-025-22171-3}{Tariq et al. (2025)} claim to analyze factors affecting student stress using \href{https://www.kaggle.com/datasets/rxnach/student-stress-factors-a-comprehensive-analysis/data}{a Kaggle dataset} on the subject. Both articles reproduce a substantial portion of the original dataset uploader's \href{https://github.com/Chhabii/Students-Stress-Factors-Analysis}{sample analysis}, as detailed by \href{https://pubpeer.com/publications/2DDE0CBFB7C2138A5945A5292EF624#2}{Maarten Van Kampen on PubPeer}.

It is important to remember that Kaggle is only a content-sharing platform. No dataset should be trusted merely because it is on Kaggle. Kaggle hosts many valuable datasets uploaded by trustworthy entities, but it also hosts unverifiable datasets unsuitable for scientific research.

\href{https://doi.org/10.64898/2026.02.24.26347028}{Gibson et al. (2026)} present evidence that some commonly-analyzed Kaggle datasets lack data provenance and might be fabricated. In late 2025, Springer Nature \href{https://doi.org/10.53053/XFKN9077}{retracted} nearly 40 publications based on a dataset supposedly containing the faces of autistic and non-autistic children. This dataset, created by one retired engineer, contains no evidence that its subjects are or are not autistic as claimed nor that they provided consent to be included. It is not possible to make credible claims based on analyzing a dataset like this because there is no reason to trust its provenance.

Note that Kaggle provides a 10-point ``usability'' score for datasets uploaded to its platform. This score cannot be used to identify if datasets are real or trustworthy.

\subsubsection*{Additional resources}

\begin{itemize}
    \setlength\itemsep{-0.5em}
    \item \href{https://doi.org/10.64898/2026.02.24.26347028}{``Evidence of Unreliable Data and Poor Data Provenance in Clinical Prediction Model Research and Clinical Practice'' (2026)}
    \item \href{https://doi.org/10.53053/XFKN9077}{``Exclusive: Springer Nature retracts, removes nearly 40 publications that trained neural networks on `bonkers' dataset'' (2025)}
\end{itemize}

\end{document}